% =====================================================================
%  PART I -- THE REDUCTION
%  (merged edition v3, section 1)
% =====================================================================
\clearpage
\part{The reduction}
\label{part:reduction}

This part does one thing: it writes the mass gap as an exact statement about a single number, and
then shows why that rewriting does not move the problem. The rewriting is exact; the obstruction is
that each of its three conjuncts, when one tries to prove it, requires itself. We state the
equivalence (\S\ref{sec:reduction-equiv}), identify the one framework in which the object can be
placed and show that everything that framework gives for free is a \emph{transport} of positivity
rather than a production of it (\S\ref{sec:reduction-framework}), take the three legs one at a time
and exhibit the circularity of each (\S\ref{sec:reduction-legs}), record the one sentence that
summarises the situation (\S\ref{sec:reduction-nomove}), and close with the qualification that
governs the whole part: $w_1$ is a property of the operator \emph{pair}, not of the Hamiltonian
(\S\ref{sec:reduction-operator}).

% ---------------------------------------------------------------------
\section{The equivalence}
\label{sec:reduction-equiv}

Fix the object once and for all. Let $O=\tfrac14\operatorname{tr}F^2$, let $\nu\ge 0$ be its
K\"all\'en--Lehmann spectral measure on $[0,\infty)$, and put
\begin{equation}
  M_0 \;:=\; \int d\nu \;=\; \operatorname{Var}(O),\qquad
  m_1^2 \;:=\; \inf\operatorname{supp}\nu,\qquad
  f_1^2 \;:=\; \nu(\{m_1^2\}),\qquad
  w_1 \;:=\; \frac{f_1^2}{M_0}\;\in\;[0,1].
  \label{eq:w1}
\end{equation}
So $w_1$ is the atom fraction at the bottom of the spectrum: the weight of the lightest excitation
in the source $O$, normalised by the total spectral weight. The bound $w_1\in[0,1]$ is an identity,
not an estimate: $0\le \nu(\{m_1^2\})\le\nu([0,\infty))=M_0$, for every positive measure, every
$(L,a)$, and every normalisation of $M_0$
(\code{research/v3/species/U01.md}, \S3). The K\"all\'en--Lehmann representation itself is not an
assumption of this paper but a theorem of the frame: $\nu\ge0$ and $f_1^2\ge0$ follow from
reflection positivity and the spectral theorem \cite{kl,os1973}.

\smallskip
\noindent\textbf{One notation, fixed once.} The variable carried by $\nu$ is $m^2$, so the atom of
the lightest state sits at the \emph{point} $m_1^2$ and is written $\nu(\{m_1^2\})$; the correlator
is then $C(t)=\int e^{-mt}\,d\nu(m^2)=\int e^{-\sqrt{s}\,t}\,d\nu(s)$ with $s=m^2$. Writing
$\nu(\{m_1^2\})$ and $d\nu(m)$ in the same display, as an earlier draft did, is a notational
collision of exactly the kind catalogued in \S\ref{sec:ledger}; the convention above is used
throughout.

\begin{claimbox}{cC}
\claimhead{Reduction}{.} With $w_1$ as in \eqref{eq:w1}, with $\beta=2.5$ the physical coupling of
this paper, and with the physical operator $O=\tfrac14\operatorname{tr}F^2$ held fixed,
\begin{align*}
  \textbf{Clay}\ \text{\emph{on the atom route}}\;\Longleftrightarrow\;&
  \underbrace{\inf_{L,t}\, w_1(L,2.5,t) \,>\, 0}_{\text{(i) the uniform floor at }\beta=2.5}
  \;+\;\underbrace{\exists\ \text{a monotone (or convergent) step in } L}_{\text{(ii) the volume step}}\\[2pt]
  &+\;\underbrace{\exists\,\Delta_0>0\ \ \forall \beta\ge\beta_0:\ m_{0^{++}}(\beta)\ge\Delta_0}_{\text{(iii) the }\beta\text{-tail-uniform floor}} .
\end{align*}
Three qualifications belong to the statement and are not commentary on it.
\emph{(a) Route.} ``On the atom route'' names the class of arguments that certifies the gap through
the atom weight of a fixed $O$; \S\ref{sec:survey-three} exhibits two solvable models in which $w_1$
and $m$ point in opposite directions, so the equivalence fails outside that class, and a proof
outside it cannot be obtained by refining (i)--(iii).
\emph{(b) Operator.} Leg (i) is stated at fixed $O$; \S\ref{sec:reduction-operator} shows that
dropping the qualification lets the quantity sweep all of $[0,1]$ at fixed spectrum.
\emph{(c) Floor, not sign.} Leg (i) asks for an infimum, not for pointwise positivity. The two differ
by exactly the compactness that \S\ref{sec:reduction-framework} shows one does not have, and it is
the pointwise form alone that the measurement of \S\ref{sec:measurement} supplies. The three
conjuncts are not three routes to one target: they are the three things a proof of the target must
supply, and the equivalence is a rewriting of the quantifiers, not a theorem about Yang--Mills.
\end{claimbox}
\src{Reduction and the three-leg split: \code{research/v3/species/U15.md}, \S1 and \S5 (the target
must be written as $w_1(\beta)\ge w_0(\beta)$ with $w_0(2.5)$ unknown); \code{U02.md}, \S5(c) (the
target is a full-family infimum, $\inf_{L,a,t}w_1\ge w_0$, not an endpoint value; at fixed $\beta$
the family is the $\inf_{L,t}$ written in the statement, since $a=a(\beta)$);
\code{UNIF3.md}, \S0 and \S4.1.}

The one conjunct that is not merely a restatement is (iii), and it is where the equivalence earns
its keep, because (iii) turns out to be the same statement as a continuity statement in the lattice
spacing.

\begin{claimbox}{cT}
\claimhead{Lemma}{.} (Tail identity.) Suppose $\beta\mapsto a(\beta)$ is strictly decreasing with
$a(\beta)\to 0$ (asymptotic freedom). Then for every $c_0>0$,
\[
  \exists a_0\ \forall a<a_0:\ m_{0^{++}}(a)\ge c_0\Lambda
  \quad\Longleftrightarrow\quad
  \exists\beta_0\ \forall\beta>\beta_0:\ m_{0^{++}}(\beta)\ge c_0\Lambda .
\]
Hence ``$a$-uniform'' and ``$\beta$-tail-uniform'' are not two requirements but one, written in two
units: $a<a_0\iff\beta>\beta(a_0)$, on the same quantity, in the same units.
\end{claimbox}
\begin{proofx}
Both sides are of the form ``$m\ge c_0\Lambda$ on a terminal set of a monotone parameter''. The map
$\beta\mapsto a(\beta)$ is a strictly monotone bijection onto its range, so it carries terminal sets
to terminal sets and preserves the inequality. No estimate on $m$ is used, and none is produced.
\end{proofx}
\src{\code{research/v3/species/UNIF3.md}, Lemma A (\S1.2); the arithmetic is the project's own, the
input ``$a\to 0\iff\beta\to\infty$'' is \code{research/uniqueness/REPORT.md}:43 as quoted there.}

The point of the lemma is negative and it is worth stating separately: it removes a whole class of
proposed second stages. There is no fifth direction called ``$a\to 0$'' that could be attacked after
the physical coupling has been settled, because a statement at a \emph{single} $\beta$ (including
$\beta=2.5$) constrains no sequence with $a\to 0$. What (iii) asks for is uniformity on a
\emph{tail}, and that is what (ii) and (i) also ask for, in their own variables.

\begin{table}[htbp]
\centering\small
\begin{tabularx}{\textwidth}{@{}l l l X@{}}
\toprule
Leg & Object & Status & Source \\
\midrule
(i)  & $\inf_{L,t}w_1(L,2.5,t)>0$ at fixed $O$ & \LO{} as a floor; \LE{} only pointwise & \code{RUN3.md}, \S3; \code{U11.md}, \S3.2 \\
(ii) & monotone/convergent step in $L$ & \LN{} for the present toolkit & \code{U02.md}, \S2; \code{U03.md}, \S2--\S3 \\
(iii) & $\beta$-tail-uniform floor $=\Delta_0$ & \LO{} (value); \LN{} (target eligibility) & \code{H1-FINAL.md}, \S0, \S3 \\
(iii)$'$ & ($a$-uniform form of (iii)) & \LT{} identical to (iii) & \code{UNIF3.md}, \S1.2 \\
\bottomrule
\end{tabularx}
\caption{The three legs, their objects, and their grades. Row (i) is the uniform floor; the pointwise
form, which is what a measurement can reach, is strictly weaker (\S\ref{sec:reduction-framework}).
The equivalence of rows (iii) and (iii)$'$ is the tail identity proved above.}
\label{tab:legs}
\end{table}

% ---------------------------------------------------------------------
\section{The unique framework}
\label{sec:reduction-framework}

Because $w_1$ is an atom mass, the object $\nu$ lives in a space where compactness costs nothing: a
variational upper bound (Ritz; the bottom of the spectrum is a minimum of a quadratic form
\cite{reed-simon}) gives $m_1^2\le\Lambda$, so $\nu/M_0$ is a probability measure on the
compact $K=[0,\Lambda]$, and $\mathcal{P}(K)$ is weak-$*$ compact and metrisable by
Banach--Alaoglu with the Prokhorov tightness hypothesis satisfied automatically. That is the whole
of the good news, and it is the reason the compactness class of arguments is out of play rather
than merely unpolished.

\begin{claimbox}{cT}
\claimhead{Theorem}{.} (Transport, not production.) Let $K$ be compact, $w_1:K\to[0,1]$ continuous,
and $w_1(x)>0$ for every $x\in K$. Then
\[
  \inf_{x\in K} w_1(x)\;=\;\min_{x\in K} w_1(x)\;=\;w_1(x_0)\;>\;0 ,
\]
and, conversely, $\inf_K w_1>0$ implies pointwise positivity on $K$. The two statements are
equivalent: compactness plus continuity adds nothing to pointwise positivity except the licence to
move the infimum to the point where it is attained.
\end{claimbox}
\begin{proofx}
A continuous function on a compact set attains its minimum (Weierstrass, one line; no Ascoli--Arzel\`a
is needed, since a single function is trivially equicontinuous). The converse is immediate. Hence
the two formulations are equivalent and the compactness hypothesis is not doing work.
\end{proofx}
\src{\code{research/v3/species/U01.md}, \S1.1--\S1.2 (``compactness does not produce positivity'' is
stated there as the whole content of the route); \code{U16.md}, \S3.2 for the choice of $K$ and the
freeness of tightness.}

So on a truncated window $K=\{L_1,\dots,L_K\}\times[a_{\min},a_0]$ with $a_{\min}>0$ the floor
$\inf_K w_1=w_1(L^*,a^*)>0$ is available, and it is an \emph{identity}: it carries a positivity that
was already assumed at every point of $K$ to a single point of $K$. It does not create positivity,
and it says nothing about $a_{\min}\downarrow 0$, because $a=0$ is not in $K$.

The second free item is the direction of semi-continuity, and its direction is the wrong one.

\begin{claimbox}{cT}
\claimhead{Proposition}{.} (Atoms are upper semi-continuous.) If $\nu_n\Rightarrow\nu$ weakly and
$x$ is fixed, then $\limsup_n \nu_n(\{x\})\le\nu(\{x\})$. Consequently, if the limit measure has no
atom at $m_1^2$, then $w_1^{(n)}\to0$: at the limit the atom can be destroyed, and upper
semi-continuity does not prevent it. A floor, by contrast, needs a \emph{lower}
semi-continuous bound.
\end{claimbox}
\begin{proofx}
By the Portmanteau theorem applied to the closed sets $\overline{B_\varepsilon(x)}$,
$\nu_n(\{x\})\le\nu_n(\overline{B_\varepsilon(x)})$, so
$\limsup_n\nu_n(\{x\})\le\nu(\overline{B_\varepsilon(x)})$; let $\varepsilon\downarrow0$.
\end{proofx}
\src{\code{research/v3/species/U01.md}, \S4.3, Proposition P-2; \code{U16.md}, \S0(c) and \S3.3.}

Weak-$*$ semi-continuity is therefore \emph{free in the direction ``finite volume $\to$ limit''} and
absent in the direction a floor requires, namely ``limit $\to$ nearby members''. The gap is not
technical: there is a compact family, with a limit point that carries an atom of mass $0.3$, whose
uniform floor is exactly $0$.

\begin{table}[htbp]
\centering\small
\begin{tabular}{@{}rrrrrr@{}}
\toprule
$h$ & mass & $\int s\,d\mu_h$ & $E=\text{bottom atom}$ & $\inf\operatorname{supp}$ & $\|\mu_h-\mu_0\|_{TV}$ \\
\midrule
$0.000$ & $1.00000$ & $2.40000$ & $0.300000$ & $1.00000$ & $0.000$ \\
$0.200$ & $1.00000$ & $2.40000$ & $0.000000$ & $0.80000$ & $0.600$ \\
$0.100$ & $1.00000$ & $2.40000$ & $0.000000$ & $0.90000$ & $0.600$ \\
$0.050$ & $1.00000$ & $2.40000$ & $0.000000$ & $0.95000$ & $0.600$ \\
$0.020$ & $1.00000$ & $2.40000$ & $0.000000$ & $0.98000$ & $0.600$ \\
$0.005$ & $1.00000$ & $2.40000$ & $0.000000$ & $0.99500$ & $0.600$ \\
$0.001$ & $1.00000$ & $2.40000$ & $0.000000$ & $0.99900$ & $0.600$ \\
\bottomrule
\end{tabular}
\caption{A compact family whose member at $h=0$ carries an atom of mass $0.3$ and whose uniform atom
floor is $0$. Total variation is quoted on the $L^1$ scale,
$\|\mu-\nu\|_{TV}=\int|d\mu-d\nu|$, so that the value $0.600$ is twice the usual normalisation
$\frac12\int|d\mu-d\nu|=0.300$; the constant offset is a convention, the vanishing of the floor is
not. Computed (\code{U16.md}, \S3.3).}
\label{tab:compact}
\end{table}

\noindent The family is $\mu_h=0.3\,\mathrm{Unif}[1-h,1+h]+0.7\,\delta_3$ on $\mathcal{P}([0,3])$,
with $\mu_0=0.3\,\delta_1+0.7\,\delta_3$. Every moment readout is constant to six digits
($\int s\,d\mu_h\equiv2.400000$), the family is compact, the limit point carries an atom of mass
$0.3$, and yet $\inf_{h\in[0,1/2]}E(\mu_h)=0$ while $\sup=0.3$ is attained only at $h=0$. Compact
family plus a limit atom does \emph{not} give a uniform atom floor (\code{U16.md}, \S3.3, machine
computed).

Finally, neither end of the parameter space is compact, and this is what removes the tame
theorems. At $a=0$ and at $L=\infty$ the family has \emph{no member}: the domain is open in $a$ and
the $L$'s are discrete with no interpolation (\code{U01.md}, \S2.3). The dimension grows like
$(L/a)^4\to\infty$ (\code{DF-CONTINUITY.md}:247, as quoted in \code{U15.md}), so no o-minimal or
other tameness hypothesis of finite dimension applies, and a uniformity theorem for a definable
family cannot be invoked even in principle. The honest formulation is that the obstruction is not
``not compact'' but ``not closed'': the family has no member at the point where its infimum is
sought.

\begin{claimbox}{cN}
\claimhead{No-go}{.} (The generator slot is empty.) The seven general tools for uniformity over an
infinite-dimensional family --- Banach--Steinhaus, Ascoli--Arzel\`a, Prokhorov, Baire category,
$\Gamma$-convergence, uniform large deviations, and induction on $L$ --- are all transports or
promotions: pointwise $\to$ uniform, local $\to$ global, weak $\to$ strong, $L\to L+1$. Not one of
them turns a non-uniform family into a uniform one. Banach--Steinhaus is the only one that imposes no
structure on the index set, and its hypothesis is literally ``the family is already uniform at every
point''; applied to $w_1$ it returns the trivial bound $1$.
\end{claimbox}
\src{\code{research/v3/species/U16.md}, \S1 (statements), \S2 (the table of what supplies
uniformity), \S4 (the $a\to0$ side is outside the range of all seven).}

% ---------------------------------------------------------------------
\section{Why each leg is circular}
\label{sec:reduction-legs}

\subsection{(i) Positivity needs the atom, and the atom is what is missing}

The leg asks for $w_1>0$, and $w_1$ is an atom weight. By the proposition above, the atom functional
is upper semi-continuous only, while a floor requires lower semi-continuity, and the failure is
realised: a thin atom $p_n\delta_{m_1^2}+(1-p_n)\delta_\Lambda$ with $p_n\to0$ has $w_1=p_n\to0$ with
$M_0$, $m_1^2$ and the gap all fixed (\code{Q6-H1-STATUS.md}:208--215), and there are gapped
atomless measures with $f_1^2=0$ and every moment, positivity and Hankel condition satisfied
(\code{R6-ATOM-ADVERSARIAL.md}:254--259). So compactness together with pointwise positivity does not
imply a uniform floor at the level of measures, and the disproof is by explicit witnesses rather
than by a failed search (\code{U01.md}, Proposition P-3).

The instruments do not repair this, and the reason is structural rather than quantitative. Every
instrument built in this campaign is a weak-$*$ continuous functional of $\nu$: the moments
$\int s^k\,d\nu$, the Ursell combination \code{$\phi^{U}$}, and the coefficient sequence
\code{$a_n$}. They are blind to atoms in a precise sense. For any weak-$*$
continuous $g$ one has $\sup_\nu|E(\nu)-g(\nu)|\ge f_1^2/2$, i.e.\ the error of the instrument is at
least the signal (\code{R3-ATOM-FUNCTIONAL.md}, Theorem R3-2, as recorded in \code{U15.md}, \S3);
with $M_0$ and the first moment fixed, $w_1$ ranges over an interval
(\code{WA-ALIGNMENT.md}:37--39, recorded in \code{U02.md}, \S6); with $(M_0,\chi)$ fixed, $f_1^2=1$
and $f_1^2=0.25$ are both
realisable (\code{LF-LASTFOUR.md}:221); and all of them --- $\operatorname{Var}$, Hessian, $E[H]$,
$\chi$, $f_1^2$ itself as a readout --- are translation-invariant normalised moments, hence
structurally blind to position (\code{LOOP-RUNBOOK.md}, \S69.1, the \code{tb01} verdict). The
variance families used to build them are the Efron--Stein ones \cite{es}, and the one lower-bound
mechanism of a different type, a transport--entropy criterion for a spectral gap \cite{vrs}, needs a
curvature input that the lattice measure at $\beta=2.5$ does not supply. Fourteen
named routes were checked against this target and none closes it: $0/14$, eleven no-go, two
conditional with the hypothesis equal to the conclusion, one numerical
(\code{research/v3/species/U15.md}, \S2). That is the circularity of leg (i): the positivity must be
produced by an instrument that can see an atom, and every instrument available is continuous in the
topology in which atoms are destroyed.

\subsection{(ii) Monotonicity is not there, and would not suffice}

The leg asks for a monotone step in the volume. There is no monotonicity theorem for $w_1(L)$, and no
$L$-sequence of measurements either; within the known families both directions occur. The
delocalised lowest mode gives $w_1=2/(L-1)$, i.e.\ $2,1,\tfrac23,\tfrac12,\tfrac25,\tfrac13,\tfrac27$
for $L=2,\dots,8$, strictly decreasing to $0$ (exact rational check, \code{U02.md}, \S2.2); a
localised light mode with finite correlation length gives a constant. The volume dependence of a
mass is in general a genuine $1/L$ expansion rather than a monotonicity \cite{luscher77}, which is
one more reason to expect no sign here. So $w_1(L)$ is not a property of the species, and no sign can
be extracted from the structure alone.

Even granting monotonicity, it does not give the leg, because a non-increasing sequence can decrease
to zero: $2/(L-1)\searrow0$, so ``non-increasing'' does not imply ``infimum positive''. And the one
positivity input available does not even supply monotonicity in the right family: reflection
positivity yields $m_{\mathrm{eff}}(t)\downarrow$, $\widehat{w}_1(T)\downarrow$ with
$\widehat{w}_1(T)\ge w_1$, and $\lambda_2^t\downarrow$ with $\lambda_2\le1$ --- all in the transfer
direction, all upper bounds on $w_1$, hence the opposite direction to a floor; in the spatial
direction reflection positivity yields existence, not monotonicity, and $w_1(L)$ is not a member of
its monotone family at all (\code{U03.md}, \S2--\S3). The leg is therefore not merely unproved: the
kind of object that would prove it does not exist in the toolkit.

\subsection{(iii) The tail floor is \texorpdfstring{$\Delta_0$}{Delta\_0}, and \texorpdfstring{$\Delta_0$}{Delta\_0} is dead}

The leg asks for a floor uniform on the coupling tail. That object is $\Delta_0$, and $\Delta_0$ has
already been judged: it is neither sufficient nor necessary for the gap, so it cannot serve as a
form of the target, and its value is open.

\begin{claimbox}{cN}
\claimhead{No-go}{.} ($\Delta_0$ is not a rewriting of the gap.) Two implications fail.
\emph{Insufficiency}: a thin atom has $\Delta_{\mathrm{dyn}}>0$ while $m=0$
(\code{Q6-H1-STATUS.md}:53--56, 110--114; machine-checked as \code{H1\_does\_not\_imply\_atom} with
no \code{sorry}). \emph{Unnecessity}: in two-dimensional Yang--Mills $\Delta_{\mathrm{dyn}}=0$ while
$m=3g^2/8>0$ (\code{N03-WFULL.md}, \S0; \code{H1-FINAL.md}, \S3.3). Hence
$\Delta_0>0\not\Rightarrow m>0$ and $m>0\not\Rightarrow\Delta_0>0$: the floor is neither a
sufficient nor a necessary form of the target.
\end{claimbox}
\src{\code{research/v3/species/H1-FINAL.md}, \S0 rows 4--5 and \S3.3; \code{N04-2D.md}:14,49 for
$G''=W_{\mathrm{geo}}+\Delta_{\mathrm{dyn}}$ with $W_{\mathrm{geo}}\ge2$ and
$\Delta_{\mathrm{dyn}}=O(e^{-2t})\to0$ in two dimensions.}

\smallskip
\noindent\textbf{Where the missing floor actually sits.} The operator whose spectral gap is wanted is
the Hessian of the action at a background, $G''=-D^{2}+2adF_B$, and its indefinite direction at a
constant chromomagnetic field is not an artefact of the expansion: the Savvidy--Nielsen--Olesen
background \cite{savvidy,no} carries a negative mode in the transverse-magnetic channel
\cite{leutwyler}, and the instability is physical rather than a gauge copy. This is why the coarse-mode
gap cannot be had from positivity of $G''$: $G''$ is not positive, so the measure that would supply
the floor sits partly on the negative side. What is wanted is therefore not a sharper estimate of
$G''$ but a statement about the spectrum of $\nu$ that does not route through $G''$ at all --- which
is leg (iii) again.

The solvable model $M6$ is where the insufficiency used to be argued, and that argument must now be
quoted with its quantifier attached: on $M6$ one has $\Delta_{\mathrm{dyn}}=2\omega$ on the
$a$-mode while $m_6=\min(\omega,2\pi^2/L^2)\to0$, but the $b$-mode gives $\Delta_{\mathrm{dyn}}=0$,
so $\inf_v\Delta_{\mathrm{dyn}}=0$ and $M6$ does \emph{not} satisfy the target under the
$\forall v$ reading of the quantifier (\code{H1-FINAL.md}, \S3.2, correcting
\code{N12-VERDICT.md}). The insufficiency witness is the thin atom; the unnecessity witness is two
dimensions.

What remains is a single open cell, and it is not a gap in an argument: it is the absence of a
$(L,a)$-uniform positive lower bound on $\Delta_{\mathrm{dyn}}$, equivalently a uniform spectral gap
of the coarse-mode operator $G''$ (\code{H1-FINAL.md}, \S0; \code{N11-FALLBACK.md}, \S6). That cell
is exactly the content of leg (iii), which is why the leg is circular: to prove it one must already
have it.

% ---------------------------------------------------------------------
\section{Why it does not move}
\label{sec:reduction-nomove}

\begin{qbox}
\textbf{The equivalence is itself the reason it does not move.} Each of the three legs, the moment
one sets out to prove it, turns out to require itself: (i) needs an atom seen by an instrument that
cannot see atoms, (ii) needs a monotone quantity that the positivity in hand does not supply and
that would not suffice if it did, and (iii) \emph{is} the missing cell, renamed.
\end{qbox}

This is worth separating from the ordinary experience of a hard problem, where a reduction shortens
the distance to the target. Here the reduction is exact and the distance is unchanged: the three
conjuncts are not three lemmas that would compose into a proof, but three descriptions of one
missing object, and the descriptions are equivalent precisely because they all name the same
uniformity. A reduction that is an equivalence of this kind cannot be a source of progress; it can
only be a source of bookkeeping, and the bookkeeping is what the rest of this paper performs.

% ---------------------------------------------------------------------
\section{The key qualification: \texorpdfstring{$w_1$}{w1} depends on the operator}
\label{sec:reduction-operator}

The last point is the one that changes how leg (i) must be written, and it is a two-line
counterexample rather than an empirical worry.

\begin{claimbox}{cT}
\claimhead{Theorem}{.} ($w_1$ is not a function of the spectrum.) Fix any Hamiltonian $H$ with
bottom $E_0$, first excitation $\psi_1$, second excitation $\psi_2$ and $\lambda_1=E_1-E_0>0$. For
$\theta\in[0,\pi/2]$ and $0<\varepsilon\ll1$ set
\[
  O_\theta \;=\; 1+\varepsilon\bigl(\cos\theta\,\psi_1+\sin\theta\,\psi_2\bigr),
\]
which is a legitimate non-negative source. Then $\operatorname{Var}(O_\theta)=\varepsilon^2$ and
$f_1^2=\varepsilon^2\cos^2\theta$, so
\[
  w_1(O_\theta)\;=\;\cos^2\theta\;\in\;[0,1],
\]
while $\operatorname{Spec}(H)$, $\lambda_1$, $L$ and $a$ are unchanged. Hence there is no function $F$
with $w_1=F(\operatorname{Spec}H)$; $w_1=f_1^2/M_0$ is a property of the pair $(H,O)$, not of $H$.
\end{claimbox}
\begin{proofx}
Insert $O_\theta$ into \eqref{eq:w1}. The variance is $\varepsilon^2$ by orthonormality of
$\Omega,\psi_1,\psi_2$; the bottom atom is the $\psi_1$ component squared, $\varepsilon^2\cos^2\theta$;
the ratio is $\cos^2\theta$. Nothing in the spectrum enters the computation.
\end{proofx}
\src{\code{U11.md}, \S3.2, equations (3.2)--(3.3), with the scope note there: this is a no-go on the
operator-pair class, not a claim that $w_1$ is spectrum-independent on Yang--Mills, where $(H,O)$ is a
single fixed pair (file: \code{research/v3/species/U11.md}).}

The counterexample is not a pathology of an unphysical source; it is realised in the measurement.
On one Markov chain, with the same $12000$ configurations, two operators give atom fractions that
differ by a factor $20.2$. That a mass estimate is an estimate \emph{for the operator that was
interpolated}, not for the theory, is the working premise of lattice spectroscopy
\cite{morningstar}; what the theorem above adds is that the dependence is total --- it can exhaust
the whole interval $[0,1]$ --- so no amount of statistics can remove it.

\begin{table}[htbp]
\centering\small
\begin{tabularx}{\textwidth}{@{}l r r r r r@{}}
\toprule
basis & $C(0)$ & $C(1)/C(0)$ & $m_{\mathrm{fit}}$ & $w_1$ (mass pinned) & $S/N(2)$ \\
\midrule
\code{raw}   & $19.888$ & $0.07372\pm0.00265$ & $1.707$ & $0.3427\pm0.0119$ & $5.0$ \\
\code{TAPE1} & $45.527$ & $0.00302\pm0.00297$ & $1.073$ & $0.0168\pm0.0142$ & $0.9$ \\
\code{TAPE2} & $39.428$ & $0.00150\pm0.00298$ & $0.762$ & $0.0107\pm0.0136$ & $1.6$ \\
\bottomrule
\end{tabularx}
\caption{One Markov chain, one set of configurations, three operators, all read at $m_1=1.546$
(Table~\ref{tab:m-bases}). Only the operator changes, and the atom fraction changes by a factor $20.2$.}
\label{tab:opdep-front}
\end{table}

\noindent Geometry $8^3\times12$, $\beta=2.5$, $\kappa=1.25$, $N=12000$ ($8$ chains $\times$ $1500$
configurations, $\mathrm{gap}=3$, $400$ thermalisation sweeps per chain), plaquette
$0.65234\pm0.00002$, errors from a $20$-block delete-one jackknife, $w_1$ with the mass pinned at
$m_1=1.546$. \code{TAPE} is a spatial APE smearing with the time direction left unsmeared
(\code{RUN3.md}, \S2--\S3). The two leading entries differ by the factor
$0.3427/0.0168=20.2$ reported in the abstract; the \code{raw} value reproduces an earlier
independent run at $0.02\sigma$.

Two consequences follow, and the second is a boundary of this paper.

\begin{v3box}
\textbf{(1) Leg (i) must be stated at fixed $O$.} ``$w_1>0$'' has no content until $O$ is named: at
fixed spectrum the value ranges over $[0,1]$. The correct reading of (i) is
\emph{for the physical operator $O=\tfrac14\operatorname{tr}F^2$, $w_1(2.5)>0$ at each $(L,a)$}, and
every statement in this paper is to be read that way.

\textbf{(2) The measured value is \LE{} and nothing more.} It is a lattice measurement at
$L_t=12$, $L_s=6$ and $L_s=8$, at one coupling, with $S/N(2)=5.0<10$: the $28.8\sigma$ quoted in the
abstract is significance from zero, not a demonstration of a floor, and the deficit in $S/N$ is
recorded rather than repaired (\code{RUN3.md}, \S4).
\end{v3box}

\begin{warnbox}
\textbf{Where the grades must not be guessed.} Three hooks in this part are grade-sensitive and are
flagged here rather than settled.

(a) The $M6$ row: the insufficiency claim is true per mode ($\Delta_{\mathrm{dyn}}=2\omega$ on the
$a$-mode) but false under the $\forall v$ quantifier ($\inf_v=0$), and the repository withdrew the
row on that ground (\code{H1-FINAL.md}:102,138). This part therefore uses the thin atom, not $M6$, as
the insufficiency witness; any sentence restoring $M6$ to that role would need a grade that is not
available in the sources.

(b) The factor $20.2$ compares two \emph{operators}, not two theories: it is evidence for the
operator dependence, and it is not evidence for or against any floor.

(c) Leg (i) at fixed $O$ and the mass gap as usually stated differ by an $O$-independence
requirement, and no such statement exists in this paper. That is a genuine boundary of the
reduction, not a technicality: the reduction is exact at fixed $O$, and the passage to a statement
about $H$ alone is exactly what the counterexample forbids.
\end{warnbox}

% ---------------------------------------------------------------------
\section{Where the reduction sits in the literature}
\label{sec:reduction-inherited}

Almost nothing in this part is new except the reduction itself, and a reader is entitled to know which
sentences carry an obligation and which are borrowed. The following inventory is organised by what is
borrowed, not by chronology; every item is used somewhere in the paper, and none of them is used at a
grade higher than the source supports.

\noindent\textbf{(a) The frame.} The problem is the one posed in \cite{jaffe}: a quantum field theory
satisfying the Wightman axioms \cite{sw}, with a mass gap above the vacuum. The spectral data
$\nu$, $m_1$, $f_1^2$ of \eqref{eq:w1} are the K\"all\'en--Lehmann data of that frame \cite{kl}, and
the positivity that makes $\nu$ a measure at finite volume is reflection positivity
\cite{os1973}, in the lattice form of Osterwalder--Seiler \cite{osterwalder-seiler}.

\noindent\textbf{(b) The strong-coupling mass gap is not ours.} That lattice gauge theory has a mass
gap at large $\beta$ is a theorem of the 1970s and 1980s: Seiler's lecture notes \cite{seiler}, the
Brydges--Giffen--Durhuus--Fr\"ohlich polymer expansion \cite{bgdf}, and, for the lower-dimensional
models, Glimm--Jaffe--Spencer \cite{gjs} and Feldman--Osterwalder \cite{fo}. The criterion that the
five rings of \S\ref{sec:est-i1} are run against is Koteck\'y--Preiss \cite{kp}; the edge $1/48$ of
the strong-coupling window is the threshold of Shen--Zhu--Zhu's log-Sobolev and Poincar\'e
inequalities \cite{szz}. None of
these is reproved here. What \S\ref{sec:est-i1} adds is the residue statement $I1$ and, with it, the
removal of a premise that the earlier literature needed --- that the $0^{++}$ channel be strictly the
lightest.

\noindent\textbf{(c) The renormalisation-group route.} The programme that would prove the gap at all
couplings by a multiscale expansion is Balaban's \cite{balaban,balaban-conv}, and its completion with an infrared
cutoff is Magnen--Rivasseau--S\'enor \cite{mrs}. Neither is used in this paper, and the reason is
recorded rather than hidden: the estimates are not uniform in the ways \S\ref{sec:survey} needs, and
the constants in them are not the constants this campaign can control.

\noindent\textbf{(d) Analytic tools.} Log-concavity enters through Pr\'ekopa \cite{prekopa} and
Brascamp--Lieb \cite{bl}; the spectral statements through Kato \cite{kato} and Reed--Simon
\cite{reed-simon}; level crossings through von Neumann--Wigner \cite{vnw}; the operator-comparison
bounds of row 12 of \S\ref{sec:survey-table} through Temple \cite{temple}; the stochastic and
semigroup side through Feller \cite{feller}, Jerrum--Sinclair \cite{js}, Sinclair \cite{sinclair},
Bauerschmidt--Brydges--Slade \cite{bbs} and the lecture notes of \cite{chatterjee}.

\noindent\textbf{(e) Other routes to a mass scale, and why none of them is the missing machine.}
Lattice formulations and their numerical study are Wilson's \cite{wilson} and Creutz's \cite{creutz},
with the step-scaling and flow technologies of \cite{luscher-stepscaling,luscher-flow} and the
spectroscopy of \cite{morningstar}; the large-$N$ and loop-equation route is
\cite{thooft-planar,eguchi-kawai,makeenko-migdal}; the confining-string route is
\cite{polchinski-strominger,luscher-weisz}; the holographic route is
\cite{maldacena,gkp,witten,witten-thermal,polchinski-strassler,sakai-sugimoto}; the conformal-bootstrap
route is \cite{elshowk,prv}; the resurgent route is \cite{abs}; the $\beta$-function and
$a$-theorem complex is \cite{caswell,vrvl,banks-zaks,polchinski,komargodski-schwimmer,dksst,zamolodchikov,osborn},
with the conformal-collider bounds of \cite{hofman-maldacena,hofman-li} and the near-conformal lattice
studies of \cite{aoki}; and the infrared scenarios that put the mass scale in the gauge sector are
Gribov's \cite{gribov}, its topological formulation by Singer \cite{singer} and Zwanziger
\cite{zwanziger}, the Dyson--Schwinger truncations of \cite{vsmekal,fischer-alkofer} reviewed in
\cite{huber}, and the dynamically generated gluon mass of \cite{cornwall}. Confinement itself has
been proved in a restricted model, three-dimensional $U(1)$ \cite{gopfert}.

Each entry in (e) supplies an upper bound, a scenario, or a numerical value at a fixed coupling. None
of them supplies the object \S\ref{sec:verdict-needed} names, and the uniform direction of
\S\ref{sec:survey} is the observation that this is not an accident of any one route.

\label{part:reduction-end}
